\chapter{Research Methodology and Experimental Design}\label{chap:method}
\sloppy

This chapter describes the experimental design. The project is an experimental, comparative study organised around four tracks: Plain BP, Secure BP, Plain FF and Secure FF. It started from plaintext baselines, used bounded secure \ac{bp} prototypes to establish feasibility under \ac{mpc}, and has since completed one full-size secure \ac{bp} run that is matched to the plaintext \ac{bp} baseline. It is designed to continue with a secure \ac{ff} implementation, a controlled comparison, per-operation profiling and one \ac{mpc}-aware modification of \ac{ff}. The parts that have been carried out are described as such; the remaining parts are described as planned.

\section{Study Design and Experimental Tracks}\label{sec:method-design}
The independent variable of the main study is the learning mechanism (\ac{bp}, \ac{ff} and, later, one \ac{mpc}-aware \ac{ff} variant). A second factor is the execution mode: plaintext floating-point training in PyTorch, or secure fixed-point training in MP-SPDZ. Crossing the two gives the four tracks of Figure~\ref{fig:method-tracks}. The dependent variables are utility (loss and accuracy on held-out data) and cost (runtime, communication volume, communication rounds where measurable and, later, a per-operation breakdown).

\begin{figure}[ht]
\centering
\begin{tikzpicture}[
  track/.style={draw, rounded corners, align=center, minimum width=5.4cm, minimum height=2.0cm, font=\small},
  done/.style={track, very thick, fill=black!7},
  partial/.style={track, very thick},
  todo/.style={track, dashed},
  link/.style={-{Stealth[length=2.5mm]}, thick},
  biglink/.style={{Stealth[length=2.5mm]}-{Stealth[length=2.5mm]}, thick},
  lab/.style={font=\footnotesize, align=center}
]
\node[done] (pbp) {\textbf{Plain BP}\\PyTorch, CPU, float32\\completed: 5 seeds (42--46)};
\node[done, right=2.6cm of pbp] (pff) {\textbf{Plain FF}\\PyTorch, CPU, float32\\completed: 5 seeds (42--46)};
\node[partial, below=2.3cm of pbp] (sbp) {\textbf{Secure BP}\\MP-SPDZ, Rep3 Ring, 3 parties\\full-size: 1 run (seed 42)\\bounded: 3 historical runs};
\node[todo, below=2.3cm of pff] (sff) {\textbf{Secure FF}\\MP-SPDZ, Rep3 Ring, 3 parties\\not implemented};
\draw[biglink] (pbp) -- node[lab, above]{plaintext\\benchmark} (pff);
\draw[link] (pbp) -- node[lab, left]{paired\\comparison\\(seed 42)} (sbp);
\draw[link, dashed] (pff) -- node[lab, right]{planned} (sff);
\draw[biglink, dashed] (sbp) -- node[lab, below]{RQ1--RQ3\\(pending)} (sff);
\end{tikzpicture}
\caption{The four experimental tracks and their status on 16 November 2026. Solid boxes: completed evidence; thick unfilled box: partial evidence; dashed: not yet implemented.}
\label{fig:method-tracks}
\end{figure}

\section{Learning Rule versus Secure-Computation Protocol}\label{sec:method-separation}
Two layers of the system are kept apart throughout the thesis. The \emph{learning rule} (\ac{bp} or \ac{ff}) defines which quantities are computed and how parameters are updated. The \emph{secure-computation protocol} (here, MP-SPDZ's three-party replicated-ring protocol) defines how each arithmetic or nonlinear operation is evaluated on secret-shared values, and therefore what each operation costs. A change at one layer is not a change at the other. In particular, rewriting secure \ac{bp} from scalar loops into matrix products (Chapter~\ref{chap:impl}) changes the cost of evaluating \ac{bp} under the same protocol, but the program still implements \ac{bp}: the same forward pass, the same softmax cross-entropy loss, the same backward recursion and the same \ac{sgd} update. Conversely, the research question about \ac{ff} concerns a change of learning rule under a fixed protocol. Any comparison of \ac{bp} and \ac{ff} under \ac{mpc} must therefore hold the protocol, threat model, arithmetic and implementation style (for example, matrix formulation) constant, or the effect of the learning rule cannot be separated from the effect of the implementation.

\section{Comparisons and Controlled Variables}\label{sec:method-comparisons}
The design calls for several distinct comparisons. They control different variables and support different claims, and they are kept separate in the results. Table~\ref{tab:method-comparisons} lists them with what is matched and what differs.

\begin{table}[ht]
\centering
\small
\begin{tabular}{@{}>{\raggedright\arraybackslash}p{2.6cm}>{\raggedright\arraybackslash}p{4.4cm}>{\raggedright\arraybackslash}p{4.3cm}>{\raggedright\arraybackslash}p{2.4cm}@{}}
\toprule
Comparison & Matched & Differs & Status \\
\midrule
C1: Plain BP vs Plain FF & Data split, seeds 42--46, hidden widths, batch size, number of epochs per trained unit, CPU software stack & Learning rule; parameters (118{,}282 vs 116{,}992); learning rate; optimizer updates (200 vs 400); presentations; inference procedure & Completed (5 seeds) \\
C2: Plain BP vs Secure BP (seed 42) & Split, sample order, initial parameters, architecture, loss, learning rate, batch schedule, epochs, evaluated checkpoint (epoch 5) & Arithmetic (float32 vs fixed point with probabilistic truncation); execution environment; nonlinear approximations & Completed (1 seed) \\
C3: Secure BP vs Secure FF & Planned: framework, protocol, threat model, arithmetic, data, host, implementation style & Learning rule and its budget & Not started \\
C4: Secure FF vs modified Secure FF & Planned: everything except the modification & Modification (with ablation) & Not started \\
Historical: A0, A1, BP-COMP-001 & 784--8--10 network, 90/30/30 subset, protocol & Ordering, initialisation, ReLU derivative, capture coverage & Completed; not compared causally \\
\bottomrule
\end{tabular}
\caption{Comparisons in the study, the variables each one controls, and their status.}
\label{tab:method-comparisons}
\end{table}

Comparison C2 deliberately crosses execution modes: it asks whether the secure implementation reproduces the plaintext learning trajectory for the same seed. It is a correctness and fidelity check, not a statement that secure training is equivalent to plaintext training in general. Hardware is not one of the controlled variables of C2: the plaintext benchmark records its software stack but not its processor, and the secure run records only that all parties ran on the local host (Section~\ref{sec:setup-env}). Runtime is therefore not compared across execution modes. For the planned comparison C3, the design calls for at least five seeds for utility and three repetitions for timing. The secure results available today are single runs and do not yet satisfy this requirement.

\section{Algorithm-Specific Training Budgets}\label{sec:method-budget}
\ac{bp} and \ac{ff} cannot be trained with identical procedures, so ``matched settings'' cannot mean identical work. In the plaintext benchmark, \ac{bp} trains the whole 784--128--128--10 network end to end for five epochs: 40 mini-batches per epoch (39 of 128 examples and one of 8) give $5\times40=200$ optimizer updates and $5\times5{,}000=25{,}000$ example presentations. \ac{ff} trains a 784--128--128 network greedily, one layer at a time: layer~1 is trained for five epochs, then frozen, and layer~2 is trained for five epochs on the normalised output of layer~1. Each mini-batch is presented once with the correct label (positive) and once with a uniformly random wrong label (negative), so \ac{ff} performs $2\times5\times40=400$ optimizer updates and $2\times2\times5\times5{,}000=100{,}000$ presentations, and every presentation of a training example to layer~2 also requires a forward pass through the frozen layer~1. \ac{ff} uses learning rate $0.01$, \ac{bp} uses $0.05$. At inference, \ac{bp} needs one forward pass per example, whereas \ac{ff} needs one pass per candidate label (ten passes) and chooses the label with the highest goodness accumulated over both layers. These differences are documented rather than equalised; differences in accuracy or time between the plaintext methods are therefore not attributed to the learning rule alone, and the budget of each method must be reported again with the secure comparison.

\section{Protocol, Threat Model and Fixed-Point Configuration}\label{sec:method-protocol}
All secure experiments use MP-SPDZ v0.4.4 \cite{keller2020mpspdz} at commit \texttt{113c4e47\allowbreak dd94bca9\allowbreak ecab029d\allowbreak 2f0abe64\allowbreak 9bea5c52}, with the three-party, semi-honest, honest-majority replicated-ring protocol (Rep3 Ring) over a 64-bit ring: at most one of the three parties may be corrupted, and a corrupted party is assumed to follow the protocol. Fixed-point numbers use \texttt{sfix(16,31)}: 16 fractional bits and 31 bits in total, giving a nominal range of about $[-16384,16384)$ and a quantisation step of $2^{-16}$ (Section~\ref{sec:bg-fixed}). For this protocol the MP-SPDZ compilation path enables probabilistic truncation and a three-way share split (\texttt{use\_trunc\_pr} and \texttt{use\_split(3)}); these are defaults of the execution path and not separately tuned choices. The compiler output lists 7{,}424{,}425 integer probabilistic truncations for the bounded runs A1 and BP-COMP-001 and 91{,}680{,}035 for the full-size run. The log of the historical run A0 instead ends with MP-SPDZ's suggestion to enable these two options, which indicates that A0 was executed without them. The deterministic alternative \texttt{sfix.round\_nearest} was not used, and no claim of zero leakage is made.

All secure runs since A1 also print the virtual-machine message that the computation domain is too small for low-round truncation and would need at least 103 bits. A reading of the pinned source (Section~\ref{sec:disc-trunc}) indicates that the message reports which truncation variant could not be used; its numerical and security implications for the executed variant have not been analysed and are listed as an open limitation. A separate malicious-security protocol (MASCOT) was explored in a tutorial during environment setup but is not the protocol of any reported experiment. Security statements apply only to this configuration.

\section{Historical Bounded Secure BP Study}\label{sec:method-scenario}
The first secure \ac{bp} experiments were deliberately small so that end-to-end secure runs were feasible with a scalar implementation. They are retained as development and validation milestones and are not combined with the full-size experiment.

\paragraph{Scenario.} A 784--8--10 fully connected network (6{,}370 trainable parameters, ReLU hidden layer, linear output layer) is trained on a 90/30/30 train/validation/test subset of MNIST. Three parties each own 30 training, 10 validation and 10 test examples; ownership is assigned by sample index modulo three, with three training examples per class and one validation and one test example per class for each party. Each mini-batch has size three, one example from each party, which gives 30 mini-batches per epoch; three epochs make 90 stochastic-gradient updates at learning rate $0.05$ without momentum or weight decay. Softmax cross-entropy is computed in a numerically stable form, and only aggregate results (losses and correct counts) are released.

\paragraph{Ordering of training examples.} An audit of the historical runs A0 and A1 found that their batches were class-major: each batch contained three examples of the same class, and the same fixed schedule was repeated in every epoch, with no example skipped or duplicated. Long runs of one class can bias \ac{sgd}; the effect on the historical runs has not been isolated. For the comparison-readiness run BP-COMP-001, a deterministic per-party shuffle of the training order in every epoch was introduced. It preserves the selected sample indices, the party ownership, the 90/30/30 split and the one-example-per-party-per-batch structure, and it is reproducible from an ordering seed (42042).

\paragraph{Parameter initialisation.} The historical runs used a public, deterministic initialisation of the weights provided by one party, $W_{ij}=(i-j)/1000$, with zero biases; this is neither private nor random. Candidate schemes were compared in controlled plaintext diagnostics C0--C3, and configuration C3 was selected for BP-COMP-001. C3 combines per-party shuffled ordering, He initialisation of the hidden layer \cite{he2015delving}, Xavier (Glorot) uniform initialisation of the output layer \cite{glorot2010understanding} and zero biases, with initialisation seed 42: the 784$\times$8 hidden weights are drawn from a normal distribution with variance $2/784$ and the $8\times10$ output weights uniformly from $[-\ell,\ell]$ with $\ell=\sqrt{6/(8+10)}$. The detailed C0--C2 results are not part of the archived artifacts, so no statement is made that C3 outperformed every alternative. \emph{Ordering and initialisation were changed together in C3}, so any difference between the historical runs and BP-COMP-001 cannot be attributed to either factor alone.

\section{Full-Size Matched Secure BP Experiment}\label{sec:method-full}
The full-size experiment (identifier SECURE-BP-MATCHED-v1, batched-matrix program) is designed to reproduce the Plain BP baseline configuration under \ac{mpc} as closely as the arithmetic allows, so that comparison C2 is a paired comparison for one seed.

\paragraph{Data and split.} MNIST is split with the same function as the plaintext baseline: 5{,}000 training and 1{,}000 validation examples from a seeded permutation of the official training set, and 1{,}000 test examples from positions 1{,}000--1{,}999 of a seeded permutation of the official test set (positions 0--999 are excluded because they had been exposed in early development).

\paragraph{Order, batches and initialisation.} The training order is the order produced by the plaintext baseline's seeded PyTorch data loader (batch size 128, shuffling, seed 42), with the loader state continuing across the five epochs. Each epoch has 39 batches of 128 examples and one final batch of 8, so five epochs give 200 updates. The initial parameters are the PyTorch default initialisation of the same model after setting seed 42, transposed from PyTorch's output-by-input layout to the input-by-output layout of the secure program. The input bundle records SHA-256 hashes of the five epoch orders, and an independent verifier checked that the bundle reproduces the baseline order, ownership and parameter layout (Section~\ref{sec:impl-provenance}).

\paragraph{Parties and disclosure.} Party~0 inputs the initial parameters. Every training, validation and test example is input by the party whose index equals the example's position modulo three, so each party owns about one third of the data. Features, labels, activations, gradients and parameters remain secret-shared throughout training; only the mean training loss of each epoch and the aggregate validation and test loss and correct counts are opened. The initial parameters are derived from a public seed and are therefore not private to party~0; the experiment protects the training data, not the initialisation.

\paragraph{Learning rule.} Training is plain \ac{bp} with mini-batch \ac{sgd} at learning rate $0.05$, no momentum and no weight decay; the gradient of each batch is averaged over the actual batch size, so the final batch is divided by 8, as in PyTorch.

\section{Evaluation Policy}\label{sec:method-eval}
The plaintext benchmark selects, for every seed, the checkpoint with the highest validation accuracy over the training epochs and reports its test accuracy. The secure program evaluates only the final (epoch-5) model. For seed 42, the plaintext baseline selected epoch 5, so comparison C2 compares the same checkpoint. For the five-seed plaintext aggregate this is not the case: seed 44 of Plain BP selected epoch 3, and every Plain FF seed selected the first epoch of layer~2. Five-seed plaintext aggregates are therefore not directly comparable with a future multi-seed secure aggregate unless the same checkpoint policy is used; this choice must be fixed before further secure seeds are run.

\section{Verification and Measurement Definitions}\label{sec:method-measure}
Correctness of the secure programs is checked at several levels: by an aggregate verifier for the bounded runs; by a strict ReLU-boundary program on public values; by matched plaintext FP64 references (for BP-COMP-001, for a single batch of the matrix prototype, and for the one-epoch smoke run of the full-size program); and, for the full-size run, by the paired comparison C2 with the recorded Plain BP run. These comparisons test close agreement of aggregate or intermediate quantities; they do not establish bit-exact equivalence or identical predictions.

Cost is measured with the following definitions, applied per experiment:
\begin{itemize}
  \item \emph{MP-SPDZ time}: the \texttt{Time} value printed by party~0. The run output states that it includes preprocessing (offline phase). It is started after the program has been loaded and does not include compilation.
  \item \emph{External elapsed time}: the wall-clock time measured by \texttt{/usr/bin/time -v} around the command that compiles the program and launches the three parties; available for the full-size and smoke runs only.
  \item \emph{Party-0 data sent} and \emph{party-0 rounds}: the volume sent by party~0 and its approximate count of communication rounds, as printed by the run.
  \item \emph{Global data sent}: the sum of the data sent by all three parties, as printed by the run.
  \item \emph{Compiler VM rounds}: the virtual-machine round count printed by the compiler; a compile-time statistic.
  \item \emph{Maximum resident set size}: as printed by \texttt{/usr/bin/time -v} for the timed process tree.
\end{itemize}
Runs whose measurement boundaries or capture coverage differ are not compared.

\section{Planned Secure FF Methodology}\label{sec:method-planned}
The design continues with a secure \ac{ff} implementation on the same data split, protocol, threat model, fixed-point type, host and matrix implementation style as the full-size secure \ac{bp} program, starting from the Plain FF configuration (784--128--128, label embedded in the first ten input pixels, L2 normalisation between layers, goodness as the sum of squared ReLU activities, threshold $2.0$, softplus local loss, per-layer \ac{sgd}, uniformly random wrong labels, inference over all candidate labels). Design questions that must be settled before implementation include how the wrong label of a negative example is generated without revealing the true label, how the softplus loss and the L2 normalisation are evaluated (exactly or by approximation), how inference over ten candidate labels is batched, and which checkpoint is evaluated. After a matched comparison with repeated runs (C3), both secure methods will be profiled per operation, and one bottleneck-motivated \ac{mpc}-aware modification of \ac{ff} will be evaluated against the unmodified baseline with an ablation (C4). None of these steps has been carried out at the time of writing, and no result of them is reported in this thesis.
